\section{Introduction: Why the Photon Mass Matters}


The question of whether the photon is exactly massless is not merely 
philosophical. It probes the foundations of electrodynamics, the internal 
consistency of quantum field theory, and the empirical limits of our 
ability to test nature at large scales. Classical electromagnetism assumes 
a massless photon from the outset, yet this assumption should not remain 
unquestioned: even an extremely small, non-zero photon mass would have 
detectable physical consequences.

A finite photon mass would modify Maxwell’s equations, alter the structure 
of electromagnetic fields, and violate the exact long-range behaviour of 
Coulomb interactions. The associated Proca theory predicts a finite 
interaction range, a longitudinal polarization mode, and distinct 
propagation properties of electromagnetic radiation. None of these effects 
have been observed, but their theoretical implications provide valuable 
guidance for establishing quantitative limits.

Modern physics offers several independent paths to constrain the photon 
mass. Astrophysical observations probe electromagnetic fields over 
galactic and interplanetary scales, yielding bounds far beyond the reach 
of laboratory experiments. Within quantum electrodynamics, the masslessness 
of the photon is intimately connected to gauge symmetry and the 
renormalizability of the theory. Even the behaviour of entangled photons 
reveals the underlying two-state helicity structure that characterizes a 
massless spin-1 particle.

This chapter outlines the motivation for examining the photon mass from a 
unified perspective. The subsequent sections combine symmetry principles, 
field-theoretic arguments, astrophysical data, precision measurements and 
the helicity structure of quantum entanglement to build a coherent case: 
if the photon has a mass at all, it must be extraordinarily small. The 
goal is not only to summarise the current limits, but to clarify why 
masslessness is a natural and stable property within the theoretical 
framework of electromagnetism.
